\section{Related Work}
\label{sec:related}

\noindent\textbf{Data Profiling.}
The closest line of work to \system\ is \emph{data profiling}, which aims to discover constraints directly from data instances~\DIFdelbegin %DIFDELCMD < \cite{profSurveyFelix2015,Huhtala1998TANE,Bohannon2007CFD,Fan2009DiscoverCFD,Papenbrock2016HyFD,Bleifuss2017Hydra,chu2013discovering,Pena2019DiscoveryDC,Pena2022FastDC,icdePena2024}%%%
\DIFdelend \DIFaddbegin \cite{profSurveyFelix2015,Huhtala1998TANE,Fan2009DiscoverCFD,Papenbrock2016HyFD,Bleifuss2017Hydra,chu2013discovering,Pena2022FastDC,icdePena2024}\DIFaddend . This includes formalisms such as functional dependencies, conditional dependencies, and denial constraints (DCs). DCs are particularly relevant in our setting because they subsume keys, FDs, and ordering constraints~\cite{chu2013discovering}, making them expressive enough to capture many cross-tuple rules in regulatory text. \DIFdelbegin \DIFdel{The key difference, however, is that profiling is driven by patterns in }\DIFdelend \DIFaddbegin \DIFadd{Approximate discovery tolerates violations in dirty data~}\cite{Pena2019DiscoveryDC,martin2026lima}\DIFadd{, user-specific criteria help experts rank discovered rules~}\cite{fan2023topk}\DIFadd{, and LLMs now propose candidate rules from schemas and samples~}\cite{biester2024llmclean,glock2026ledqer,li2026lcats}\DIFadd{. Still, all these methods are driven by }\DIFaddend the observed data. \DIFdelbegin \DIFdel{It can }\DIFdelend \DIFaddbegin \DIFadd{They }\DIFaddend reveal what the data seems to satisfy, but not whether that pattern is \DIFdelbegin \DIFdel{actually }\DIFdelend required by policy or regulation. As a result, profiling cannot recover constraints that are already violated in the released data, and \DIFdelbegin \DIFdel{recent work shows that }\DIFdelend discovered DCs often include many spurious constraints~\cite{martin2025false}. Our goal is to recover constraints from normative text and then use them to validate the data.
\altic{[R3-O1, R1-O1] Completar o related work: (i) citar e discutir descoberta APROXIMADA de DCs (VLDB 2019) e trabalho de guiar especialistas na definição de regras, alinhando com a correção feita na intro; (ii) adicionar REEs (regras com ML-predicates) como linha relacionada e contrastar o espaço de predicados deles com o nosso vocabulário controlado.}

\smallskip
\noindent\textbf{Data Cleaning.}
Data cleaning systems use constraints, statistical signals, or both to detect and repair errors~\cite{dataerrorsZiawash,chu2013holistic,holoclean,raha2019,peng2022garf,cleaningDeploy2024}\DIFaddbegin \DIFadd{, including rules with ML predicates (REEs)~}\cite{bao2024rock} \DIFadd{and LLM-based detectors~}\cite{ni2025zeroed,zhang2024cleaning}\DIFaddend . These systems complement our setting. Constraint-based methods are interpretable and auditable, but they assume that the relevant rules are already available. Data-driven methods reduce the need for manual rule specification, but without an external source of truth, they can struggle to distinguish rare valid values from genuine errors~\cite{cleaningDeploy2024}. Our work addresses this missing step of obtaining such rules from regulatory documents. We use DCs \DIFaddbegin \DIFadd{over a controlled, schema-derived predicate space }\DIFaddend as the output formalism\DIFdelbegin \DIFdel{because of practical implementation }\DIFdelend \DIFaddbegin \DIFadd{, so every constraint is executable and auditable, }\DIFaddend and, as prior work has shown, \DIFdelbegin \DIFdel{once DCs are available, they }\DIFdelend \DIFaddbegin \DIFadd{DCs }\DIFaddend can be checked efficiently at scale\DIFdelbegin \DIFdel{using specialized validation engines }\DIFdelend ~\cite{facet2021,liu2024rapidash,youri2024}.

\smallskip
\noindent\textbf{Regulatory NLP.}
A separate line of work uses NLP to analyze regulations and compliance-related documents, for example, by extracting rights and obligations or by checking whether requirement statements comply with legal or policy text~\DIFdelbegin %DIFDELCMD < \cite{breaux2006regulatory,breaux2008analyzing,kiyavitskaya2008automating,amaral2023nlp,compliance2025ACMTSEM}%%%
\DIFdelend \DIFaddbegin \cite{breaux2006regulatory,kiyavitskaya2008automating,amaral2023nlp,hassani2024rethinking,compliance2025ACMTSEM,masoudifard2024graphrag}\DIFadd{. LLMs now extract structured rules from regulations~}\cite{guliani2026dejure,hettiarachchi-etal-2025-code,kabzhan2025semantic} \DIFadd{and translate financial regulations into compliance code~}\cite{li2025compliance2code}\DIFadd{, although they still struggle to apply long sets of interdependent rules~}\cite{yu2026ruleworld,soni2026tam}\DIFaddend . These approaches remain \DIFdelbegin \DIFdel{primarily }\DIFdelend text-centered \DIFdelbegin \DIFdel{, focusing on classification, extraction, or comparison of regulatory statements}\DIFdelend \DIFaddbegin \DIFadd{or target a fixed application model}\DIFaddend . Our problem starts from similar documents, but we need to connect regulatory statements to a relational schema and translate them into constraints whose violations can be checked in the data.

\smallskip
\noindent\textbf{Text-to-SQL.}
Text-to-SQL is also related, since it maps natural-language input to a SQL query over a database schema~\DIFdelbegin %DIFDELCMD < \cite{yu2018spider,picard2021,resdsql2023,pourreza2023din,dail2025,llmTEXT2SQLSurvey2025,redsql2025,XiYanSQL}%%%
\DIFdelend \DIFaddbegin \cite{yu2018spider,picard2021,resdsql2023,pourreza2023din,dail2025,llmTEXT2SQLSurvey2025,redsql2025,XiYanSQL,alrashed2026verification,rabonato2026portuguese}\DIFaddend . The difference lies in \DIFdelbegin \DIFdel{in }\DIFdelend how the generated logic is used, as Text-to-SQL produces a query that answers a user question on the current database instance, whereas \system\ aims at declarative invariants that should hold for all valid instances of the data. In addition to mapping text to schema elements, as in Text-to-SQL, it must also identify scope conditions and define violation semantics. %DIF < Taken together, these lines of work highlight the gap addressed in this paper.
%DIF < \DIFdelbegin \DIFdel{: extracting executable data quality constraints from normative text and grounding them in a relational schema for validation}\DIFdelend \DIFaddbegin \DIFadd{Closer in output, LLMs generate OCL constraints from specifications of UML models~}\cite{abukhalaf2023ocl,attarwala2026ocl}\DIFadd{, and Doc2DB systems verify extracted relations against integrity constraints~}\cite{zhang2026doc2db}\DIFadd{; in both cases, the target schema is designed together with the constraints, whereas }\system\DIFadd{\ grounds regulatory text in an independently published dataset}\DIFaddend .
